\section{Section 1}


\subsection{Some Definitions}


\begin{question}
    Definieren hier die f-Divergenz allgemein für positive lineare Funktionale
    über von Neumann Algebren.
    Haben wir dadurch später Vorteile, oder willst du das einfach erstmal
    so allgemein wie möglich machen?
    Für die Physik und viele Resultate werden wir schätze ich mal Zustände
    verwenden.
\end{question}



\begin{definition}\label{def:f_div}
Let $\varphi, \psi:\mathcal{M}\rightarrow\mathbb{C}$ be positive lineare functionals over a von
Neumann algebra $\mathcal{M}$ such that
$\varphi \ll \gg \psi$ \gext{was das nochmal genau???}
and $f:(0,\infty)\rightarrow \mathbb{R}$ be a convex function with
$f(1) = 0$ \gext{(brauche ich das hier???)}.
The \textit{f-divergence} of $\varphi$ and $\psi$ is given by
\begin{equation}
\begin{split}
    D_f(\varphi| \psi) &= \int_{1}^{\infty}
    f''(\gamma) E_\gamma(\varphi| \psi)
    + \gamma^{-3} f''(\gamma^{-1}) E_\gamma(\psi|\varphi) d\gamma\\
    &= \int_{1}^{\infty} f''(\gamma) E_\gamma(\varphi|\psi) d\gamma
    + \int_{0}^{1} \gamma f''(\gamma) E_{\frac{1}{\gamma}}(\psi|\varphi) d\gamma.
\end{split}
\end{equation}
\end{definition}

\begin{definition}
    Let $\varphi$ and $\psi$ be as above. (sicher? oder doch Zustände?)
    The hockeystick divergence $E_\gamma(\varphi|\psi)$ is defined as
    follows:
    \begin{equation}
        E_\gamma(\varphi|\psi) := (\psi - \gamma\varphi)_+(\mathbb{1})
    \end{equation}
\end{definition}

Some facts about the hockeystick divergence.

\begin{proposition}
    Let $\varphi_1, \varphi_2$ be positive lineare functionals over a von Neumann
    algebra, and $p_1, p_2$ are 
    \gext{(pos. lin. functionals or states?)}
    \begin{equation}
        \lim_{t\rightarrow\infty} E_t(\varphi_1| \varphi_2) = \varphi_1(p_2^\perp).
    \end{equation}
    \begin{equation}
        \lim_{t\rightarrow\infty} E_t(\varphi_1, \varphi_2) = 0\quad
        \Longleftrightarrow\quad p_1\leq p_2.
    \end{equation}
\end{proposition}



\subsection{First Inequalitie}

In \cite{Sason_2016} there are a lot of inequalities for f-divergences in the
context of probability measures.

We want to generalize them for states over arbitrary von Neumann algebras
(even vNA of type III).




We start with inequalitie (84) from \cite{Sason_2016}.

\begin{question}
    Wann ist $D_f(\varphi|\psi)$ endlich?
    Evtl wenn $\psi$ normal ist?
\end{question}

\begin{proposition}
    Let $f,g: (0,\infty)\rightarrow\mathbb{R}$ convex functions
    with $f(1) = g(1)$, $f'(1) = g'(1)$
    and $f(t)\leq g(t)$ for all $t\in(0,\infty)$.
    Also let $\varphi, \psi$ be positive lineare functionals over a
    von Neumann algebra such that $\varphi\ll\gg\psi$
    \gext{($D_f(\varphi|\psi)), D_g(\varphi|\psi)$ are finite}
    \begin{equation}
        D_f(\varphi| \psi) \leq D_g(\varphi|\psi).
    \end{equation}
\end{proposition}

\begin{remark}
    Alexandrovs theorem states that every convex function defined on an
    open domain is twice differentiable almost everywhere.
\end{remark}

\begin{remark}
    By assumption $f'(1) = g'(1)$ indirectly require, that $f$ and $g$ are
    differentiable at $x=1$.
    Therefore, we cannot put in e.g. the \textit{total variational distance}
    (see \cite{Sason_2016} eq. (55)).
\end{remark}


\begin{proof}
    We begin by reformulating the f-divergence given in \cref{def:f_div}.
\begin{equation}
\begin{split}
        D_f(\varphi| \psi) &= \int_{1}^{\infty}
    f''(\gamma) E_\gamma(\varphi| \psi)
    + \gamma^{-3} f''(\gamma^{-1}) E_\gamma(\psi|\varphi) d\gamma\\
    &= \int_{1}^{\infty} f''(\gamma) E_\gamma(\varphi|\psi) d\gamma
    + \int_{0}^{1} \gamma f''(\gamma) E_{\frac{1}{\gamma}}(\psi|\varphi) d\gamma\\
    &\oben{=} \int_{1}^{\infty} f''(\gamma) E_\gamma(\varphi|\psi) d\gamma
    +\int_{0}^{1}\gamma f''(\gamma)(\frac{1}{\gamma} E_\gamma(\varphi|\psi)
    + \psi(\mathbb{1}) - \frac{1}{\gamma}\varphi(\mathbb{1}))d\gamma\\
    &= \int_{1}^{\infty} f''(\gamma) E_\gamma(\varphi|\psi) d\gamma
    +\int_{0}^{1} f''(\gamma) E_\gamma(\varphi|\psi)d\gamma
    +\int_{0}^{1} f''(\gamma) (\gamma \psi(\mathbb{1}) - \varphi(\mathbb{1}))d\gamma\\
    &= \int_{0}^{\infty} f''(\gamma) E_\gamma(\varphi|\psi)d\gamma
    + \psi(\mathbb{1}) \int_{0}^{1} \gamma f''(\gamma) d\gamma
    - \varphi(\mathbb{1}) \int_{0}^{1} f''(\gamma) d\gamma\\
    &\oben{=} [f'(\gamma) E_\gamma(\varphi|\psi)]_0^\infty
    - \int_{0}^{\infty} f'(\gamma) E_\gamma'(\varphi|\psi)d\gamma
    + \psi(\mathbb{1}) \left([\gamma f'(\gamma)]_0^1
    -\int_{0}^{1} f'(\gamma) d\gamma\right)-\varphi(\mathbb{1}) [f'(\gamma)]_0^1\\
    &\oben{=} [f'(\gamma) E_\gamma(\varphi|\psi)]_0^\infty
    -\left( [f(\gamma) E_\gamma'(\varphi|\psi)]_0^\infty
    -\int_{0}^{\infty} f(\gamma) E_\gamma''(\varphi|\psi) d\gamma
    \right)\\
    &\quad +\psi(\mathbb{1}) \left( [\gamma f'(\gamma)]_0^1
    -[f(\gamma)]_0^1 \right)
    -\varphi(\mathbb{1}) [f'(\gamma)]_0^1\\
    &= \lim_{\gamma\rightarrow\infty} f'(\gamma) E_\gamma(\varphi|\psi)
    -\lim_{\gamma\rightarrow 0} f'(\gamma) E_\gamma(\varphi|\psi)
    -\left(\lim_{\gamma\rightarrow\infty} f(\gamma) E_\gamma'(\varphi|\psi)
    -\lim_{\gamma\rightarrow 0} f(\gamma) E_\gamma'(\varphi|\psi)\right)\\
    &\quad +\int_{0}^{\infty} f(\gamma) E_\gamma''(\varphi|\psi)d\gamma
    +\psi(\mathbb{1})
    \left(f'(1) - \lim_{\gamma\rightarrow 0} \gamma f'(\gamma)
    - f(1) + \lim_{\gamma\rightarrow 0}f(\gamma)\right)\\
    &\quad-\varphi(\mathbb{1})\left( f'(1) - \lim_{\gamma\rightarrow 0} f'(\gamma)\right)\\
    &=\lim_{\gamma\rightarrow\infty}
    \left( f'(\gamma) E_\gamma(\varphi|\psi) - f(\gamma) E_\gamma'(\varphi|\psi)\right)
    +\int_{0}^{\infty} f(\gamma) E_\gamma''(\varphi|\psi)d\gamma\\
    &\quad +\lim_{\gamma\rightarrow 0} f(\gamma) (E_\gamma'(\varphi|\psi) + \psi(\mathbb{1}))
    + \lim_{\gamma\rightarrow 0} f'(\gamma) (\varphi(\mathbb{1})
    - \gamma\psi(\mathbb{1}) -E_\gamma(\varphi|\psi))\\
    &\quad + \underbrace{\psi(\mathbb{1}) f'(1) - \psi(\mathbb{1}) f(1)
    - \varphi(\mathbb{1}) f'(1)}_{=:K}
\end{split}
\end{equation}
(a) state exchange, (b), (c) partial integration
K is a constant here, fällt weg bei $D_g -D_f$, da $f(1) = g(1)$ und
$f'(1) = g'(1)$.
\bigskip

Next we want to show under which conditions $D_g(\varphi|\psi)\leq D_f(\varphi|\psi)$
holds.
\begin{equation}
\begin{split}
    D_g(\varphi|\psi) - D_f(\varphi|\psi)
    &=\int_{0}^{\infty} \underbrace{(g(\gamma)-f(\gamma))}_{\geq 0} d\gamma
    \underbrace{E_\gamma''(\varphi|\psi)}_{\geq 0}
    + \lim_{\gamma\rightarrow 0} \underbrace{(g(\gamma) - f(\gamma))}_{\geq 0}
    \underbrace{(E_\gamma'(\varphi|\psi) + \psi(\mathbb{1}))}_{\geq 0\: (a)}\\
    &\quad +\lim_{\gamma\rightarrow \infty} ((g'(\gamma)-f'(\gamma)) E_\gamma(\varphi|\psi)
    - (g(\gamma) - f(\gamma)) E_\gamma'(\varphi|\psi))\\
    &\quad +\lim_{\gamma\rightarrow 0} (g'(\gamma) - f'(\gamma))
    \underbrace{(\varphi(\mathbb{1})
    - \gamma\psi(\mathbb{1}) -E_\gamma(\varphi|\psi))}_{\leq 0\: (b)}
\end{split}
\end{equation}
(a) follow from lipschitz-continuity of $E_\gamma$;
(b) follows from strict positivity of $E_\gamma$.
\end{proof}